\section{The Grammar of First-Order Logic}

Last week, we explored this ins and outs of first-order logic informally. Now
we'll start to formalize those ideas, starting with the basic grammar of FOL.

\begin{definition}[Expressions of First-Order Logic]: We can difine an
  \label{def:fol-expression}
  \emph{expression} of FOL as any sequence of its six kinds of basic symbols:\\
  \begin{enumerate}
    \item \textbf{Predicates:} $A,B,C,\dots,=$
    \item \textbf{Names:} $a,b,c,\dots$
    \item \textbf{Variables:} $x,y,z,\dots$
    \item \textbf{Connectives:} $\enot,\eand,\eor,\eif,\eiff$ 
    \item \textbf{Quantifiers:} $\eforall,\eexists$
    \item \textbf{Brackets:} (,)
  \end{enumerate}
  The alphabetic symbols (names, predicates and variables) can be used with
  or without numeric subscripts, for example $A_1, A_2, A_3, \dots$ or
  $x_1, x_2, x_3, \dots$. \\
  \\
  In the metalanguage, we will refer to arbitary expressions with greek lower
  case letters, preferring \ealpha for terms, \ebeta for predicates and \egamma
  otherwise.
\end{definition}

We assume these basic expressions have roughly the meaning discussed informally.
But, of course, that is not important here. All we want to do with a grammar is
to say which expressions are sentences and which are not.

The most basic basic expressions in FOL are terms. 

\begin{definition}[Term]
  \label{def:fol-term}
  A \emph{term} is any name or variable.
\end{definition}

Unlike TFL, FOL needs to define its atomic sentences, since they are not mere
truth-bearers, but have constituent parts. We will refer to the sentences of
FOL as \emph{formulas}.

\begin{definition}[Atomic Formula]
  \label{def:fol-atomic-formula}
  An \emph{atomic formula} of first-order logic is defined as:\\
  \begin{enumerate}
  \item Any sentence letter (i.e. any zero-place predicate) is an atomic
    formula.
  \item If \ebeta is an n-place predicate and $\ealpha_1\dots\ealpha_n$ are
    terms, then $\ebeta(\ealpha_1,\dots,\ealpha_n)$ is an atomic formula.
  \item If $\ealpha_1$ and $\ealpha_2$ are terms then $\ealpha_1=\ealpha_2$ is an
    atomic formula.
  \item Nothing else is an atomic formula. 
  \end{enumerate}
\end{definition}

\begin{definition}[Formula]
  \label{def:fol-formula}
  A \emph{formula} of first-order logic is defined recusively as:\\
  \begin{enumerate}
    \item Any atomic formula of FOL is a formula of FOL
    \item If \ephi is a formula of FOL, then \enot\ephi is a formula of FOL.
    \item If \ephi and \epsi are formulas of FOL, then $\ephi\eand\epsi$ is a
      formula of FOL.
    \item If \ephi and \epsi are formulas of FOL, then $\ephi\eor\epsi$ is a
      formula of FOL.
    \item If \ephi and \epsi are formulas of FOL, then $\ephi\eif\epsi$ is a
      formula of FOL.
    \item If \ephi and \epsi are formulas of FOL, then $\ephi\eiff\epsi$ is a
      formula of FOL.
    \item If \ephi is a formula of FOL and \ealpha is a variable, then
      $\eforall\ealpha\ephi$ is a formula of FOL.
    \item Nothing else is a formula of FOL.
  \end{enumerate}
\end{definition}

We can recycle some of our ideas about operators and scope from TFL.

\begin{definition}[Main Logical Operator]
  \label{def:fol-main-operator}
  The main operator in a formula of first-order logic is the one that was
  introduced last, given the recursive application of the construction rules
  introduced in Def \ref{def:fol-formula}.
\end{definition}

\begin{definition}[Scope]
  \label{def:fol-scope}
  The \emph{scope} of an operator in first-order logic is the subformula for
  which it serves as the main logical operator.
\end{definition}

These definitions work exaclty as they do in TFL, except for the addition of
rules for \eforall and \eexists, which are analogous to \enot in their scopal
behavior. In Def \ref{def:fol-formula}, the scope of both
\eforall\ealpha and \eexists\ealpha is \ephi.

These definitions are important because they tell us how to characterize the
many variables that may occur in an FOL formula as \emph{free} or \emph{bound}.

\begin{definition}[Bound and Free Variables]
  \label{def:fol-bound-variable}
  In first-order logic, a variable \ealpha is \emph{bound} iff it occurs in
  the scope of a quantifier that uses it: \eforall\ealpha or \eexists\ealpha.
  Otherwise it is \emph{free}.
\end{definition}

We can now identify sentences and a special kind of formula, where all variables
are bound.

\begin{definition}[Sentence]
  \label{def:fol-sentence}
  In first-order logic, a formula \ephi is a \emph{sentence} iff it contains
  no free variables.
\end{definition}

We will adopt some bracketing simplification for sanity. For example,
$((A\eand B)\eand C)$ can be written $(A\eand B\eand C)$ and
$((A\eor B)\eor C)$ can be written $(A\eor B\eor C)$, since it will not
effect the meaning of the sentence to do so. We also allow the use of square
brackets `[' and `]' as substitutes for parens `(' and `)', to make pairs of
brackets easier to spot.

Now that we have a grammar to work with, we can see how FOL allows us to be
precise about some of the ambiguities of language that we just cannot
distinguish in a strictly sentence-level logic like TFL.
\cite[Chapter 29]{Mag25} works through some examples that are
particularly illustrative.

\begin{enumerate}
\item Everything that glitters is not gold.
\item Only young cats are playful.
\item Everyone went to see a movie.
\item The king of France is not bald.
\end{enumerate}

The first sentence has two readings, the basic \emph{every A is B} one and
another that more close reflects its idiomatic sense, something like
\emph{it is not the case that everything that glitters is gold}. 

\begin{enumerate}
\item $\eforall x(F(x)\eif\enot G(x))$
\item $\enot\eforall x(F(x)\eif G(x))$
\end{enumerate}

Note how the difference involves the relative scopes of negation and the
universal quantifier.

Similarly, for the second sentence, we get three distinct readings, normally
distinguished in ordinary English with patterns of word stress.

\begin{enumerate}
\item If an animal is playful, then it is a young cat.
\item If a cat is playful, then it is young.
\item If an animal is young anf playful, then it is a cat. 
\end{enumerate}

And these can all be easily distinguished via symbolization in FOL, again by
taking account of the relative scopes of operators and quantifiers.

\begin{enumerate}
\item $\eforall x(P(x)\eif(Y(x)\eand C(x))$ 
\item $\eforall x((C(x)\eand P(x))\eif Y(x))$ 
\item $\eforall x((Y(x)\eand P(x))\eif C(x))$ 
\end{enumerate}

In the third sentence, we see how the quantifiers interact. 

\begin{enumerate}
\item There is a specific movie and everyone saw that movie.
\item Everyone went to some movie or other.
\end{enumerate}

The two readings differ with respect to which quantifier is placed within the
scope of the other.

\begin{enumerate}
\item $\eexists x(M(x)\eand\eforall y(P(y)\eif\S(y,x)))$
\item $\eforall y(P(y)\eif\forall x(M(x)\eand S(y,x)))$
\end{enumerate}

Finally, Russell's treatment of definite desriptions, while it avoids
truth-value gaps, makes the fourth sentence ambiguous in a way that
directly impacts its truth value.

\begin{enumerate}
\item There is a king of France and he is not bald. (False)
\item There is no king of France who is bald. (True)
\end{enumerate}

And, once again, the differences come down to which operator or quantifier
takes \emph{wide scope}. 

\begin{enumerate}
\item $\eexists x(K(x)\eand\forall y(K(y)\eif x\neq y)\eand\enot B(x))$
\item $\enot\eexists x(K(x)\eand\forall y(K(y)\eif x\neq y)\eand B(x))$
\end{enumerate}

Now that we know how to make them, next week we will discuss how to interpret
sentences of first-order logic.
